\documentclass[12pt]{article}
\usepackage[ukrainian]{babel}
\usepackage{fontspec}
% --- НАЛАШТУВАННЯ ШРИФТІВ ---
\setmainfont{Schoolbook-Regular.otf}[
    Path = ../,
    BoldFont = Schoolbook-Bold.otf,
    ItalicFont = Schoolbook-Italic.otf,
    BoldItalicFont = Schoolbook-BoldItalic.otf
]
\usepackage[italic]{mathastext}
\usepackage[a4paper,margin=1.8cm,bottom=2cm,top=2cm]{geometry}
\usepackage{amsmath,amssymb}
\usepackage{tikz}
\usepackage{pgfplots}
\pgfplotsset{compat=1.16}
\usetikzlibrary{calc,patterns,angles,quotes,intersections}
\usepackage{xcolor,array,fancyhdr,enumitem,multicol}
\usepackage{colortbl}
\usepackage{tcolorbox}
\tcbuselibrary{skins,breakable}
\usepackage[hidelinks]{hyperref}
\usepackage[firstpage=false, text=@pvtr2525, scale=0.6, color=gray!12]{draftwatermark}

\binoppenalty=10000
\relpenalty=10000

\definecolor{mainGreen}{RGB}{34, 120, 64}
\definecolor{yearOrange}{RGB}{220, 100, 30}
\definecolor{titleBg}{RGB}{225, 240, 220}
\definecolor{instrBg}{RGB}{255, 240, 220}
\definecolor{instrBorder}{RGB}{220, 150, 80}

\pagestyle{fancy}
\fancyhf{}
\renewcommand{\headrulewidth}{0.4pt}
\renewcommand{\footrulewidth}{0.4pt}
\fancyhead[C]{\small\color{gray!80} Збірник НМТ \textendash{} Варіант за 19 червня}
\fancyhead[R]{\small\color{mainGreen}\textbf{\href{https://t.me/pvtr2525}{@pvtr2525}}}
\fancyfoot[L]{\small\color{gray!80}\href{https://t.me/pvtr2525}{ТГ: @pvtr2525}}
\fancyfoot[C]{\small\color{gray!80}\thepage}
\fancyfoot[R]{\small\color{gray!80}НМТ 2026}

\newcommand{\answerTable}[5]{%
\par\vspace{0.15cm}
\noindent
\begin{tabular}{|*{5}{>{\centering\arraybackslash}m{3cm}|}}
\hline
\rule[-0.15cm]{0pt}{0.6cm}\textbf{А} & \rule[-0.15cm]{0pt}{0.6cm}\textbf{Б} & \rule[-0.15cm]{0pt}{0.6cm}\textbf{В} & \rule[-0.15cm]{0pt}{0.6cm}\textbf{Г} & \rule[-0.15cm]{0pt}{0.6cm}\textbf{Д} \\
\hline
\rule[-0.2cm]{0pt}{0.75cm}#1 & \rule[-0.2cm]{0pt}{0.75cm}#2 & \rule[-0.2cm]{0pt}{0.75cm}#3 & \rule[-0.2cm]{0pt}{0.75cm}#4 & \rule[-0.2cm]{0pt}{0.75cm}#5 \\
\hline
\end{tabular}
\par\vspace{0.25cm}
}

\newcommand{\answerTableTall}[5]{%
\par\vspace{0.15cm}
\noindent
\begin{tabular}{|*{5}{>{\centering\arraybackslash}m{3cm}|}}
\hline
\rule[-0.2cm]{0pt}{0.7cm}\textbf{А} & \rule[-0.2cm]{0pt}{0.7cm}\textbf{Б} & \rule[-0.2cm]{0pt}{0.7cm}\textbf{В} & \rule[-0.2cm]{0pt}{0.7cm}\textbf{Г} & \rule[-0.2cm]{0pt}{0.7cm}\textbf{Д} \\
\hline
\rule[-0.45cm]{0pt}{1.2cm}#1 & \rule[-0.45cm]{0pt}{1.2cm}#2 & \rule[-0.45cm]{0pt}{1.2cm}#3 & \rule[-0.45cm]{0pt}{1.2cm}#4 & \rule[-0.45cm]{0pt}{1.2cm}#5 \\
\hline
\end{tabular}
\par\vspace{0.25cm}
}

\newcommand{\instructionBox}[1]{%
\par\vspace{0.3cm}
\begin{tcolorbox}[colback=instrBg, colframe=instrBorder, boxrule=0.6pt, arc=2pt,
    left=8pt, right=8pt, top=4pt, bottom=4pt]
\centering\small\bfseries #1
\end{tcolorbox}
\par\vspace{0.15cm}
}

\newcommand{\sectionTitle}[1]{%
\par\vspace{0.4cm}
\begin{tcolorbox}[colback=titleBg, colframe=titleBg, boxrule=0pt, arc=8pt, halign=center, fontupper=\Large\bfseries\color{mainGreen}]
#1
\end{tcolorbox}\par\vspace{0.3cm}}
\newcommand{\chapterTitle}[1]{\clearpage\sectionTitle{#1}}

\newcommand{\nmtAnswerBox}{%
\par\vspace{0.2cm}\noindent
Відповідь:\ \ 
\begin{tabular}{@{}*{4}{|>{\centering\arraybackslash}p{0.8cm}}|@{\hspace{0.1cm}\textbf{,}\hspace{0.1cm}}*{3}{|>{\centering\arraybackslash}p{0.8cm}}|@{}}
\hline
\rule[-0.2cm]{0pt}{0.7cm} & & & & & & \\
\hline
\end{tabular}
\par\vspace{0.3cm}}

\newcommand{\matchingGrid}{%
    \begingroup
    \renewcommand{\arraystretch}{1.15}
    \setlength{\tabcolsep}{4pt}
    \footnotesize
    \begin{tabular}{r|c|c|c|c|c|}
         \multicolumn{1}{c}{} & \multicolumn{1}{c}{\textbf{А}} & \multicolumn{1}{c}{\textbf{Б}} & \multicolumn{1}{c}{\textbf{В}} & \multicolumn{1}{c}{\textbf{Г}} & \multicolumn{1}{c}{\textbf{Д}} \\ \cline{2-6}
         \textbf{1} & \phantom{X} & \phantom{X} & \phantom{X} & \phantom{X} & \phantom{X} \\ \cline{2-6}
         \textbf{2} & & & & & \\ \cline{2-6}
         \textbf{3} & & & & & \\ \cline{2-6}
    \end{tabular}
    \endgroup
}

\newcounter{zad}
\newcommand{\zadtask}[1]{\stepcounter{zad}\par\vspace{0.3cm}\noindent\textbf{\thezad.}\ \ #1\par\vspace{0.2cm}}
\newcommand{\zadnum}{\stepcounter{zad}\noindent\textbf{\thezad.}\ \ }

\begin{document}
\thispagestyle{empty}
\vspace*{1.2cm}
\begin{center}
{\Large\bfseries\color{mainGreen} РЕАЛЬНИЙ ВАРІАНТ НМТ-2026}\\[0.4cm]
{\Huge\bfseries\color{mainGreen} ВАРІАНТ ЗА 19 ЧЕРВНЯ}\\[0.6cm]
{\large Real Test Notebook з усіма геометричними рисунками}\\[1cm]
{\large\color{mainGreen}\textbf{\href{https://t.me/pvtr2525}{Телеграм: @pvtr2525}}}
\end{center}
\clearpage

\chapterTitle{ТЕСТОВИЙ ЗОШИТ (19 ЧЕРВНЯ)}
\instructionBox{Завдання 1--15 мають по п'ять варіантів відповіді, з яких лише один правильний. Виберіть правильний варіант відповіді й позначте його.}

% ===== №1 =====
\zadtask{Дівчата вдвох ліпили вареники. Тетяна зліпила втричі більше вареників, ніж Оксана. Яку частину від усіх вареників зліпила Тетяна?}
\answerTableTall{$\dfrac{1}{2}$}{$\dfrac{1}{3}$}{$\dfrac{1}{4}$}{$\dfrac{3}{4}$}{$\dfrac{2}{3}$}

% ===== №2 =====
\noindent
\begin{minipage}[c]{0.46\textwidth}
\zadnum На рисунку зображено графік зміни температури (у $^\circ$C) $T$ повітря протягом доби. Скільки всього годин за добу температура повітря знижувалася?
\end{minipage}%
\hfill
\begin{minipage}[c]{0.52\textwidth}
\centering
\begin{tikzpicture}[x=0.3cm, y=0.3cm, thick]
    \draw[step=2, gray!40, thin] (0,-4) grid (26,8);
    \draw[->] (-1,0) -- (27,0) node[right, font=\scriptsize] {$t$, год};
    \draw[->] (0,-4.5) -- (0,8.5) node[above, font=\scriptsize] {$T$, $^\circ$C};
    \node[below left, font=\scriptsize] at (0,0) {$0$};
    \node[below, font=\scriptsize] at (2,0) {$2$};
    \node[below, font=\scriptsize] at (6,0) {$6$};
    \node[below, font=\scriptsize] at (10,0) {$10$};
    \node[below, font=\scriptsize] at (16,0) {$16$};
    \node[below, font=\scriptsize] at (24,0) {$24$};
    \node[left, font=\scriptsize] at (0,2) {$2$};
    \node[left, font=\scriptsize] at (0,6) {$6$};
    \draw[magenta!80!black, thick, smooth] plot coordinates {(0,2) (2,2) (4,0) (6,-2) (8,-2) (10,0) (14,4) (16,6) (18,5) (22,2) (24,2)};
    \fill (0,2) circle (2.5pt);
    \fill (16,6) circle (2.5pt);
    \fill (24,2) circle (2.5pt);
\end{tikzpicture}
\end{minipage}
\par\vspace{0.2cm}
\answerTable{$4$}{$6$}{$8$}{$12$}{$14$}

% ===== №3 =====
\noindent
\begin{minipage}[c]{0.54\textwidth}
\zadnum Вісь Землі (пряма $OA$) утворює кут $66{,}5^\circ$ із площиною її орбіти (див. рисунок). Визначте кут між віссю Землі $OA$ та перпендикуляром $OB$ до площини її орбіти.
\end{minipage}%
\hfill
\begin{minipage}[c]{0.44\textwidth}
\centering
\begin{tikzpicture}[scale=0.82, thick, line join=round]
    % сіре тло-квадрат
    \fill[gray!22] (-2.6,-2.6) rectangle (2.6,2.6);
    % Земля (майже на весь квадрат)
    \shade[ball color=blue!18] (0,0) circle (2.4);
    \draw[gray!55, thin] (0,0) circle (2.4);
    % схематичні материки
    \begin{scope}
        \clip (0,0) circle (2.4);
        \fill[green!35!olive!60, opacity=0.55] (-1.6,0.9) .. controls (-0.9,1.5) and (-0.2,0.9) .. (-0.5,0.3) .. controls (-1.1,0.0) and (-1.9,0.3) .. (-1.6,0.9);
        \fill[green!35!olive!60, opacity=0.55] (0.3,-0.4) .. controls (1.0,-0.2) and (1.3,-1.1) .. (0.7,-1.6) .. controls (0.1,-1.4) and (-0.1,-0.8) .. (0.3,-0.4);
        \fill[green!35!olive!60, opacity=0.55] (0.6,1.4) .. controls (1.3,1.7) and (1.7,0.9) .. (1.2,0.6) .. controls (0.7,0.7) and (0.4,1.1) .. (0.6,1.4);
    \end{scope}
    % площина орбіти (горизонталь через центр)
    \draw[red!75, thick] (-3.4,0) -- (3.0,0);
    \node[right, font=\scriptsize, align=left] at (2.6,0.0) {\textit{площина}\\ \textit{орбіти}};
    % перпендикуляр OB (вертикаль, пунктир)
    \draw[dashed] (0,-2.9) -- (0,3.0);
    \node[above, font=\small] at (-0.05,3.0) {$B$};
    % вісь OA (нахилена)
    \draw (-1.15,-2.65) -- (1.15,2.65);
    \node[above, font=\small] at (1.2,2.65) {$A$};
    % центр O
    \node[below right, font=\small] at (0.05,-0.05) {$O$};
    % прямий кут між площиною орбіти і OB
    \draw (-0.32,0) -- (-0.32,0.32) -- (0,0.32);
    % кут 66,5 між віссю OA і площиною орбіти
    \draw (0.9,0) arc (0:66.5:0.9);
    \node[font=\small] at (1.5,0.55) {$66{,}5^\circ$};
    % еліпс обертання біля A
    \draw[blue!70, thick] (1.0,2.5) ellipse (0.42 and 0.14);
\end{tikzpicture}
\end{minipage}
\par\vspace{0.2cm}
\answerTable{$22{,}5^\circ$}{$23{,}5^\circ$}{$32{,}5^\circ$}{$33{,}5^\circ$}{$113{,}5^\circ$}

% ===== №4 =====
\zadtask{Розв'яжіть систему нерівностей $\begin{cases}x - 2 > -3,\\ -2x > -8.\end{cases}$}
\answerTable{$(-1; 0{,}25)$}{$(4; +\infty)$}{$(-1; 4)$}{$(0{,}25; +\infty)$}{$(-5; 4)$}

% ===== №5 =====
\zadtask{Скільки всього ребер має чотирикутна призма?}
\answerTable{$4$}{$8$}{$9$}{$12$}{$16$}

% ===== №6 =====
\zadtask{Спростіть вираз $(9a + 2b)^2 - 18ab$, якщо $a = \dfrac{1}{3}$.}
\answerTableTall{$4b^2 + 9$}{$4b^2 - 6b + 9$}{$4b^2 + 6b + 9$}{$2b^2 - 6b + 9$}{$4b^2 + 6b + 6$}

% ===== №7 =====
\zadtask{Розв'яжіть рівняння $2x + 3 = x^2$.}
\answerTable{$1;\ 4$}{$1;\ 3$}{$-3;\ 1$}{$-3;\ -1$}{$-1;\ 3$}

% ===== №8 =====
\zadtask{Графік парної функції $y = f(x)$ проходить через точку $(-3;\, 9)$. Визначте $f(3)$.}
\answerTable{$9$}{$6$}{$3$}{$-9$}{$-3$}

% ===== №9 =====
\zadtask{Які з наведених тверджень є правильними?
\vspace{0.1cm}
\begin{enumerate}[label=\textbf{\Roman*.}, align=left, leftmargin=2.6em, labelsep=0.5em, itemsep=0.08cm, topsep=0.06cm]
\item Діаметр кола, описаного навколо прямокутного трикутника, більший за гіпотенузу.
\item Центром кола, описаного навколо будь-якого трикутника, є точка перетину його медіан.
\item У рівносторонньому трикутнику центри вписаного й описаного кіл збігаються.
\end{enumerate}
\vspace{0.1cm}}
\answerTable{лише II}{лише III}{лише I та II}{лише I та III}{I, II та III}

% ===== №10 =====
\zadtask{$\cos^2 \alpha - 1 = $}
\answerTable{$\sin^2 \alpha$}{$\sin 2\alpha$}{$1$}{$\cos 2\alpha$}{$-\sin^2 \alpha$}

% ===== №11 =====
\zadtask{Ірині на день народження друзі подарували чотири книжки. Вона пообіцяла якомога швидше їх прочитати. Скільки всього в Ірини є варіантів послідовності прочитання книг?}
\answerTable{$4$}{$6$}{$12$}{$24$}{$256$}
\newpage

% ===== №12 =====
\zadtask{У прямокутній системі координат у просторі задано сферу із центром у точці $M(-9;\, 12;\, 5)$. Визначте діаметр цієї сфери, якщо вона проходить через початок координат.}
\answerTableTall{$30$}{$26$}{$10\sqrt{10}$}{$4\sqrt{21}$}{$5\sqrt{10}$}

% ===== №13 =====
\zadtask{Визначте похідну функції $y = (2x + 1)\sqrt{x}$.}
\answerTableTall{$\dfrac{3}{2\sqrt{x}}$}{$\dfrac{6x + 1}{2\sqrt{x}}$}{$\dfrac{4x + 1}{2\sqrt{x}}$}{$\dfrac{2}{\sqrt{x}}$}{$\dfrac{1}{\sqrt{x}}$}

% ===== №14 =====
\noindent
\begin{minipage}[c]{0.56\textwidth}
\zadnum На стороні $BC$ прямокутника $ABCD$ вибрано точку $K$ так, що трикутник $AKD$ прямокутний (див. рисунок). Визначте довжину сторони $AD$, якщо $AB = 6\sqrt{3}$, а $\angle KAD = 60^\circ$.
\end{minipage}%
\hfill
\begin{minipage}[c]{0.42\textwidth}
\centering
\begin{tikzpicture}[scale=0.5, thick, line join=round]
    \coordinate (A) at (0,0);
    \coordinate (B) at (0,4);
    \coordinate (C) at (10,4);
    \coordinate (D) at (10,0);
    \coordinate (K) at (2.4,4);
    \draw (A) -- (B) -- (C) -- (D) -- cycle;
    \draw (A) -- (K) -- (D);
    \pic[draw, angle radius=0.3cm] {right angle = A--K--D};
    \pic[draw, angle radius=0.55cm, pic text={\scriptsize $60^\circ$}, angle eccentricity=1.6] {angle = D--A--K};
    \node[below left] at (A) {$A$};
    \node[above left] at (B) {$B$};
    \node[above right] at (C) {$C$};
    \node[below right] at (D) {$D$};
    \node[above] at (K) {$K$};
\end{tikzpicture}
\end{minipage}
\par\vspace{0.2cm}
\answerTable{$12$}{$18$}{$12\sqrt{3}$}{$24$}{$24\sqrt{3}$}

% ===== №15 =====
\zadtask{Визначте проміжок, якому належить корінь рівняння $5^{6x + 2} = \dfrac{1}{25}$.}
\answerTable{$[-2; -1)$}{$[-1; 0)$}{$[0; 1)$}{$[1; 2)$}{$[2; +\infty)$}
\newpage
\instructionBox{У завданнях 16--18 до кожного з трьох рядків інформації, позначених цифрами, доберіть один правильний, на Вашу думку, варіант, позначений буквою.}

% ===== №16 =====
\zadtask{Доберіть до числового виразу (1--3) його значення (А--Д).\par
\vspace{0.3cm}
\noindent
\begin{minipage}[t]{0.42\textwidth}
\textit{Числовий вираз}\par\vspace{0.15cm}
\textbf{1} \quad $\dfrac{8}{5} : \dfrac{4}{15}$\\[0.35cm]
\textbf{2} \quad $\sqrt{2} \cdot \left(\sqrt{2}\right)^3$\\[0.35cm]
\textbf{3} \quad $\log_{16} 2$
\end{minipage}%
\hfill
\begin{minipage}[t]{0.3\textwidth}
\textit{Значення виразу}\par\vspace{0.15cm}
\textbf{А} \quad $-\dfrac{1}{4}$\\[0.2cm]
\textbf{Б} \quad $4$\\[0.15cm]
\textbf{В} \quad $\dfrac{1}{4}$\\[0.2cm]
\textbf{Г} \quad $6$\\[0.15cm]
\textbf{Д} \quad $\dfrac{1}{6}$
\end{minipage}%
\hfill
\begin{minipage}[t]{0.2\textwidth}
\vspace{0.4cm}
\begin{flushright}\matchingGrid\end{flushright}
\end{minipage}}

% ===== №17 =====
\zadtask{До кожного початку речення (1--3) доберіть його закінчення (А--Д) так, щоб утворилося правильне твердження.\par
\vspace{0.3cm}
\noindent
\begin{minipage}[t]{0.46\textwidth}
\textit{Початок речення}\par\vspace{0.15cm}
\textbf{1} \quad Графіки функцій $y = 3 - x$ та $y = 1 - x$\\[0.25cm]
\textbf{2} \quad Графіки функцій $y = \dfrac{1}{x}$ та $y = x^2 - 4$\\[0.25cm]
\textbf{3} \quad Графік функції $y = 3$ та графік рівняння $x^2 + y^2 = 9$
\end{minipage}%
\hfill
\begin{minipage}[t]{0.34\textwidth}
\textit{Закінчення речення}\par\vspace{0.15cm}
\textbf{А} \quad не мають спільних точок.\\[0.15cm]
\textbf{Б} \quad мають єдину спільну точку.\\[0.15cm]
\textbf{В} \quad мають лише дві спільні точки.\\[0.15cm]
\textbf{Г} \quad мають лише три спільні точки.\\[0.15cm]
\textbf{Д} \quad мають чотири спільні точки.
\end{minipage}%
\hfill
\begin{minipage}[t]{0.14\textwidth}
\vspace{0.4cm}
\begin{flushright}\matchingGrid\end{flushright}
\end{minipage}}

% ===== №18 =====
\zadtask{На рисунку зображено паралелограм $ABCD$, рівносторонній трикутник $CDK$ і прямокутну трапецію $AKNM$, що лежать в одній площині. Менша основа трапеції дорівнює стороні трикутника, її менша бічна сторона~--- стороні $BC$ паралелограма. Узгодьте геометричну фігуру (1--3) та її площу (А--Д), якщо $MN : AK = 5 : 9$ і $MN = 10$~см.
\vspace{0.3cm}
\begin{center}
\begin{tikzpicture}[scale=0.42, thick, line join=round]
    % середня горизонтальна лінія: A ... D ... K
    \coordinate (A) at (0,0);
    \coordinate (D) at (5,0);
    \coordinate (K) at (10,0);
    % паралелограм ABCD: B,C зверху
    \coordinate (B) at (2.5,4.33);
    \coordinate (C) at (7.5,4.33);
    % трапеція AKNM знизу (прямий кут при K)
    \coordinate (N) at (10,-6);
    \coordinate (M) at (3.5,-6);
    % --- заливки ---
    \filldraw[fill=orange!25, draw=black] (A) -- (B) -- (C) -- (D) -- cycle;
    \filldraw[fill=cyan!25, draw=black] (D) -- (C) -- (K) -- cycle;
    \filldraw[fill=green!20, draw=black] (A) -- (K) -- (N) -- (M) -- cycle;
    % прямий кут при K
    \draw (9.3,0) -- (9.3,-0.7) -- (10,-0.7);
    % --- мітки рівних сторін ---
    \path (B) -- (C) node[midway, sloped, font=\small] {$||$};
    \path (D) -- (C) node[midway, sloped, font=\small] {$|$};
    \path (C) -- (K) node[midway, sloped, font=\small] {$|$};
    \path (D) -- (K) node[midway, font=\small, yshift=3pt] {$|$};
    \path (K) -- (N) node[midway, sloped, font=\small] {$||$};
    \path (M) -- (N) node[midway, font=\small, yshift=3pt] {$|$};
    % --- підписи ---
    \node[left] at (A) {$A$};
    \node[above left] at (B) {$B$};
    \node[above right] at (C) {$C$};
    \node[below] at (D) {$D$};
    \node[right] at (K) {$K$};
    \node[below] at (M) {$M$};
    \node[below] at (N) {$N$};
\end{tikzpicture}
\end{center}
\vspace{0.2cm}
\noindent
\begin{minipage}[t]{0.42\textwidth}
\textit{Геометрична фігура}\par\vspace{0.15cm}
\textbf{1} \quad трикутник $CDK$\\[0.25cm]
\textbf{2} \quad паралелограм $ABCD$\\[0.25cm]
\textbf{3} \quad трапеція $AKNM$
\end{minipage}
\hfill
\begin{minipage}[t]{0.32\textwidth}
\textit{Площа фігури}\par\vspace{0.15cm}
\textbf{А} \quad $112$~см$^2$\\[0.15cm]
\textbf{Б} \quad $140$~см$^2$\\[0.15cm]
\textbf{В} \quad $25\sqrt{3}$~см$^2$\\[0.15cm]
\textbf{Г} \quad $40\sqrt{3}$~см$^2$\\[0.15cm]
\textbf{Д} \quad $60\sqrt{3}$~см$^2$\\[0.15cm]
\end{minipage}
\hfill
\begin{minipage}[t]{0.2\textwidth}
\vspace{0.4cm}
\begin{flushright}\matchingGrid\end{flushright}
\end{minipage}}
\newpage
\instructionBox{Розв'яжіть завдання 19--22. Одержані числові відповіді запишіть у спеціально відведеному місці. Відповідь записуйте лише десятковим дробом, урахувавши положення коми. Знак «мінус» записуйте перед першою цифрою числа.}

% ===== №19 =====
\zadtask{Знайдіть площу фігури, обмеженої графіком функції $y = \dfrac{x^3}{8}$, прямими $x = 2$, $x = 4$ та віссю $x$.}
\nmtAnswerBox

% ===== №20 =====
\zadtask{Василь замовив виготовлення двох рамок для своїх картин. Вартість рамки для більшої картини на $40\,\%$ вища за вартість рамки для меншої картини. Визначте вартість (у грн) рамки для більшої картини, якщо за обидві рамки Василь заплатив $1560$~грн.}
\nmtAnswerBox

% ===== №21 =====
\noindent
\begin{minipage}[c]{0.56\textwidth}
\zadnum На рисунку зображено правильну трикутну піраміду та конус, у якому радіус основи відноситься до висоти як $5 : 12$. Радіус основи конуса дорівнює радіусу кола, уписаного в основу піраміди, а висота піраміди в $\sqrt{3}$ більша за висоту конуса. Визначте об'єм (у см$^3$) піраміди, якщо об'єм конуса дорівнює $800\pi$~см$^3$.
\end{minipage}%
\hfill
\begin{minipage}[c]{0.42\textwidth}
\centering
\begin{tikzpicture}[scale=0.6, thick, line join=round]
    % правильна трикутна піраміда
    \coordinate (a) at (0,0);
    \coordinate (b) at (2.8,0.6);
    \coordinate (c) at (1.6,1.5);
    \coordinate (ap) at (1.3,5);
    \draw (a) -- (b) -- (ap) -- cycle;
    \draw (a) -- (c) -- (b);
    \draw[dashed] (a) -- (c);
    \draw[dashed] (c) -- (ap);
    % конус
    \begin{scope}[shift={(4.6,0)}]
        \draw (0,0.5) ellipse (1.5 and 0.45);
        \draw[dashed] (-1.5,0.5) -- (0,4.6);
        \draw (-1.5,0.5) .. controls (-1.5,0.1) and (1.5,0.1) .. (1.5,0.5);
        \draw (-1.5,0.5) -- (0,4.6) -- (1.5,0.5);
    \end{scope}
\end{tikzpicture}
\end{minipage}
\par\vspace{0.2cm}
\nmtAnswerBox

% ===== №22 =====
\zadtask{Визначте \textit{суму} всіх цілих значень $a$, за кожного з яких рівняння
\[
\dfrac{25\sin x - 3a + 4}{\sqrt{(x - \pi)(2\pi - x)}} = 0
\]
має принаймні один корінь.}
\nmtAnswerBox

\newpage
% ===== ВІДПОВІДІ =====
\sectionTitle{ВІДПОВІДІ ДО ВАРІАНТА (19 ЧЕРВНЯ)}

\noindent\textbf{\large Завдання 1--15 (вибір однієї відповіді)}
\par\vspace{0.2cm}
\begin{multicols}{5}
\noindent
\textbf{1.}~Г\\
\textbf{2.}~Д\\
\textbf{3.}~Б\\
\textbf{4.}~В\\
\textbf{5.}~Г\\
\textbf{6.}~В\\
\textbf{7.}~Д\\
\textbf{8.}~А\\
\textbf{9.}~Б\\
\textbf{10.}~Д\\
\textbf{11.}~Г\\
\textbf{12.}~В\\
\textbf{13.}~Б\\
\textbf{14.}~Г\\
\textbf{15.}~Б
\end{multicols}

\par\vspace{0.3cm}
\noindent\textbf{\large Завдання 16--18 (відповідність)}
\par\vspace{0.2cm}
\noindent
\textbf{16.}\quad 1--Г,\ \ 2--Б,\ \ 3--В\\[0.15cm]
\textbf{17.}\quad 1--А,\ \ 2--Г,\ \ 3--Б\\[0.15cm]
\textbf{18.}\quad 1--В,\ \ 2--Г,\ \ 3--А

\par\vspace{0.3cm}
\noindent\textbf{\large Завдання 19--22 (коротка відповідь)}
\par\vspace{0.2cm}
\begin{multicols}{4}
\noindent
\textbf{19.}~$7{,}5$\\
\textbf{20.}~$910$\\
\textbf{21.}~$7200$\\
\textbf{22.}~$-27$
\end{multicols}

\end{document}